\label{inv:C3.8}
\StatementRecord{inv:C3.8}{proved}{}{\cref{inv:A2.1,inv:A3.2,inv:C3.1,inv:C3.2,inv:C3.7}}
Retain the standard marked family of \cref{wb:thm:retained-period-deformation},
its original vectors $v_1=(1,2,-4,0)^t$, $v_2=(-1,-3,3,0)^t$,
the untranslated original cusp fan, and any
$c\in\mathcal P=\{\operatorname{Im}c<3\log(6)/\pi\}$.
There are specified smooth forms
\[
 G\in\mathcal A^{0,1}(X),\qquad Q\in\mathcal A^{1,1}(X),
 \qquad \bar\partial G=\bar\partial Q=0,
\]
whose classes satisfy
\[
 [G]=g_\gamma,\qquad
 i_b^*[Q]_{\bar\partial}=\eta,\qquad
 \exp(12[G])=f^*\mathcal O_{\mathbb P^1}(1).
\]
Here $i_b$ is any regular fibre inclusion and
$\eta=u\wedge w+6\gamma\wedge\delta$ in the original integral marking.
The formulas \eqref{inv:eq:atlas-final-forms} and
\eqref{inv:eq:atlas-final-coefficients}, with the elementary recursion
\eqref{inv:eq:atlas-weight-derivative}, give their coefficients on a
specified countable refinement of the original atlas.  Every sum is
locally finite, also at the multiple fibres, cusp double curves and
both cusp triple points.  The form $Q$ is a complex representative;
no real $\bar\partial$-closed representative can have this fibre class.
\end{proposition}

\begin{proof}
\begin{enumerate}[label=\arabic*.,leftmargin=*]
\item \emph{The actual domains and their inverse maps.}
Use the ordinary and finite covering charts in
\cref{fam:src:A.1.2,wb:thm:retained-period-deformation}, and the cusp
charts of \CompendiumRef{I}{wb:lem:toric-chart-control} and
\CompendiumRef{I}{wb:thm:candidate-cusp-geometry}.
More explicitly, write $t_1=t$, $t_2=t-1$, $q=1/t$, and retain
the fixed disjoint special discs of those constructions.  After
restricting their charts by inclusion, write their radii as
$r_1,r_2,\epsilon$, with $r_1,r_2<1/4$, and use
\[
 V_j=\{|t_j|<r_j\},\quad
 V_0=\{|q|<\epsilon\},\quad
 V_*=B\setminus\left(
 \{|t_1|\le r_1/2\}\cup\{|t_2|\le r_2/2\}
 \cup\{|q|\le\epsilon/2\}\right).
\]
Thus only $V_*\cap V_j$ are nonempty pairwise overlaps between
distinct pieces.  All three are annuli.  The finite coordinate domains
are the original torus covers over $|s_j|^{m_j}<r_j$, with
$t_j=s_j^{m_j}$ and $(m_1,m_2)=(3,4)$; their maps to $X$ are the
original free affine quotient maps.  The middle domains are the
ordinary lifted base opens over $V_*$, times $\mathbb C^2$.

For the cusp, retain the full holomorphic matrix $C(q)$ and a closed
disc $|q|\le r_0$ on which it is holomorphic.  The original radius
satisfies
\[
 0<\epsilon<\min(r_0,1),\qquad
 |\log\epsilon|>
 2\max_{|q|\le r_0}\|-2\pi\operatorname{Im}C(q)\|.
\]
These are precisely the chart and properness bounds in the cited cusp
construction and in the proof of
\cref{wb:thm:retained-period-deformation}.  For every triangle
$\sigma$ of the original periodic triangulation, use
\[
 \mathcal U_\sigma(2,\epsilon)=
 \{(z_0,z_1,z_2): |z_i|<2,\ |z_0z_1z_2|<\epsilon\},
 \qquad q=z_0z_1z_2.
 \]
The toric characters are the basis dual to the three original
height-one ray generators.  Their integral exponent matrix is
unimodular, so its inverse gives every inverse toric overlap.
These domains cover the cusp space before quotient, including its
boundary, by the cited chart-control lemma.

All these maps are local biholomorphisms to $X$.  Their overlaps are
fibre products, retaining every deck branch.  A lattice map
$\zeta\mapsto\zeta+\Pi\lambda$ has inverse
$\zeta\mapsto\zeta-\Pi\lambda$.  A finite affine map
$(s,\zeta)\mapsto(\omega s,R(s)\zeta+\Pi(\omega s)v/m)$
has inverse
\[
 s=\omega^{-1}s',\qquad
 \zeta=R(s)^{-1}\bigl(\zeta'-\Pi(s')v/m\bigr).
\]
The cusp period map has inverse
$x\mapsto q^{-B_0n}\exp(-2\pi iC(q)n)x$.
The finite collar inverse subtracts its full translating section,
and the cusp collar inverse uses the branches
$\zeta=(2\pi i)^{-1}\log x$, whose differences are the actual
integral periods.  These are refinements and inverse maps inside the
given $X$, not a replacement compactification.

\item \emph{The regular class and its sign.}
Put
\[
 \Pi=[Z\mid I_2],\quad
 Z=\begin{pmatrix}6\mu&\tau\\\beta&\mu\end{pmatrix},\quad
 T=\begin{pmatrix}\tau&6\mu\\\mu&\beta\end{pmatrix},\quad
 K=\begin{pmatrix}0&6\\1&0\end{pmatrix},\quad D=\operatorname{diag}(1,6).
\]
The exact real marking is
\[
 q_B=(r_u,r_\gamma)^t,\quad p_B=(r_w,r_\delta)^t,
 \qquad \zeta=p_B+Tq_B,
\]
with inverse $q_B=(\operatorname{Im}T)^{-1}\operatorname{Im}\zeta$,
$p_B=\operatorname{Re}\zeta-(\operatorname{Re}T)q_B$.
For $\lambda=(n_\gamma,n_u,m_w,m_\delta)$, write
$n=(n_\gamma,n_u)^t$ and define
\begin{equation}\label{inv:eq:atlas-theta-factor}
 a_\lambda(\zeta)=
 \exp\left(2\pi i\left[n^tK\zeta+\tfrac12 n^tKZn\right]\right).
\end{equation}
Symmetry of $KZ$ makes its logarithmic composition defect the integer
$n^tK m'$.  Thus these are actual line factors.

Let $Y=\operatorname{Im}T$, $S=DY^{-1}$ and
$\varphi=2\pi q_B^tDYq_B$.  The metric coefficient $e^\varphi$
obeys $h(\zeta+\Pi\lambda)|a_\lambda|^2=h(\zeta)$.
With $F=\bar\partial\partial\log h$ and $c_1=iF/(2\pi)$ its
Chern form is
\begin{equation}\label{inv:eq:atlas-semiflat}
 Q_{\rm sf}=-\frac{i}{2\pi}\partial\bar\partial\varphi
 =-\frac i2(d\zeta-dTq_B)^tS\wedge
             \overline{(d\zeta-dTq_B)}.
\end{equation}
Expanding under the exact real inverse gives
$dq_B^tD\wedge dp_B=dr_u\wedge dr_w+6dr_\gamma\wedge dr_\delta$.
In particular the fibre class is $\eta$, with its original sign.
If $y=\operatorname{Im}\tau$, $m=\operatorname{Im}\mu$,
$r=\operatorname{Im}\beta$ and $\Delta=r-6m^2/y$, then
\[
 S=\frac1{y\Delta}
 \begin{pmatrix}r&-6m\\-6m&6y\end{pmatrix},
\]
which gives exactly the previously displayed invariant fibre form.

\item \emph{The actual middle line.}
Use the original matrices $A_1,A_2$ in \cref{fam:section:A}.
Their period transformation laws and complex-linear maps are
\[
\begin{array}{c|ccc}
j&\tau_j&\mu_j&\beta_j\\\hline
1&1-1/\tau&(1-\mu)/\tau&\beta-6(\mu-1)^2/\tau+2\\
2&-1/\tau&1+\mu/\tau&\beta-6\mu^2/\tau-3,
\end{array}
\]
\[
 R_1=\begin{pmatrix}-1/\tau&0\\(1-\mu)/\tau&1\end{pmatrix},
 \qquad R_2=\begin{pmatrix}1/\tau&0\\-\mu/\tau&1\end{pmatrix},
 \qquad R_j\Pi=\Pi_jA_j.
\]
For both generators put $F_j(\zeta)=-\zeta_1^2/(2\tau)$.
The differences
\[
 \frac{\log a_{A_j\lambda}(R_j\zeta)-\log a_\lambda(\zeta)}{2\pi i}
 -F_j(\zeta+\Pi\lambda)+F_j(\zeta)
\]
on the four original lattice generators are
$(24,0,1/2,0)$ and $(-9,0,0,0)$.  Their exponential ratios are
characters, so the four values determine the full lattice comparison.
Twelfth powers kill every defect.  The line $L_*=L_\eta^{12}$ has
the linear lifts
\[
 (\zeta,v)\longmapsto
 (R_j\zeta,e^{2\pi i\,12F_j(\zeta)}v).
\]
The Gaussian sums vanish on the order-three and order-four relations.
Applying generator two and then generator one gives
$(\tau,\mu,\beta,\zeta)\mapsto(\tau+1,\mu,\beta-1,\zeta)$
and Gaussian sum zero.  These identities give the middle line on the
actual quotient, with trivial fibre multiplier for that cusp logarithm
turn.

\item \emph{Both finite affine lifts, including their phases.}
Let $b_j=v_j/m_j$.  Use the same polynomial logarithm in
\eqref{inv:eq:atlas-theta-factor} to define $a_b$ for this rational
vector.  Since $12b_j$ is integral, the actual affine map
\[
 g_j(s,\zeta)=(e^{-2\pi i/m_j}s,
                R_j\zeta+\Pi_jb_j)
\]
has the following line lift:
\begin{equation}\label{inv:eq:atlas-affine-lift}
 k_j(s,\zeta)=\exp\left(2\pi i
 \left[12F_j(\zeta)+12\frac{\log a_{b_j}(R_j\zeta)}{2\pi i}
                  +\kappa_j\right]\right),
 \quad \kappa_1=\frac{32}{3},\quad\kappa_2=\frac{81}{8}.
\end{equation}
The factor $a_{b_j}$ uses the target period $Z_j$.
Direct composition gives $g_j^{m_j}\zeta=\zeta+\Pi v_j$.
Before adding $\kappa_j$, the logarithmic difference from
$a_{v_j}^{12}$ is $-32$ for $j=1$ and $-81/2$ for $j=2$.
The stated phases cancel these differences exactly.  Hence the power
of \eqref{inv:eq:atlas-affine-lift} is the actual lattice action,
and it defines a line $L_j$ on the original free finite quotient.
The functions are holomorphic through $s=0$: the original periods
are holomorphic there, $\tau\ne0$, and every displayed exponential
is a unit.  Thus this is a line over the full multiple-fibre
neighbourhood.  Its regular fibre Chern class is $12\eta$.

\item \emph{Cusp Cartier frames and all extending units.}
Write the original full period matrix as
\[
 Z=sB_0+C(q),\quad s=\frac{\log q}{2\pi i},\quad
 B_0=\begin{pmatrix}0&1\\-1&0\end{pmatrix}.
\]
The period $n\in\mathbb Z^2$ acts by
$x\mapsto q^{B_0n}\exp(2\pi iC(q)n)x$.
Use the factor
\begin{equation}\label{inv:eq:atlas-cusp-factor}
 A_n=x^{K^tn}q^{h(n)}\exp(\pi i n^tKC(q)n),
 \qquad h(n)=-3n_\gamma^2+\frac{n_u(n_u-1)}2.
\end{equation}
Both $KB_0=\operatorname{diag}(-6,1)$ and $KC(q)$ are symmetric.
Consequently multiplication proves the exact cocycle in
\eqref{inv:eq:atlas-cusp-factor}.  Its $q$ exponent is integral.
Its ratio to \eqref{inv:eq:atlas-theta-factor}, on a fixed punctured
fibre, is $q^{-n_u/2}$, independent of $x$; its fibre Chern class
is therefore $\eta$.

Interpolate $L(i,j)=i(i-1)/2-3j^2$ on the prescribed triangulation.
Both triangles in the square with lower corner $(i,j)$ have
\[
 m_{ij}=(i,-6j-3),\quad
 l_{ij}=-i(i+1)/2+3j(j+1),\quad f_{ij}=x^{m_{ij}}q^{l_{ij}}.
\]
These are Cartier frames.  The successive ratios $f_{i+1,j}/f_{ij}$
and $f_{i,j+1}/f_{ij}$ are $x_1q^{-i-1}$ and
$x_2^{-6}q^{6(j+1)}$; their exponents vanish on the corresponding
common face, so they are units on the toric overlap.  The two triangles
inside a square have ratio one.  In the line coordinate $w=z/f_{ij}$
the overlap multiplier is the reciprocal frame ratio.

For $v=B_0n$ the exact identity
$L(u+v)-L(u)=(K^tn)\cdot u+h(n)$ cancels every $x$ and $q$ power
in the period action on $w$.  The remaining factor is
\begin{equation}\label{inv:eq:atlas-cusp-unit}
 U_{n,ij}(q)=\exp\left(\pi i n^tKC(q)n
       -2\pi i\,m_{i+v_1,j+v_2}^{t}C(q)n\right).
\end{equation}
It is a holomorphic unit at zero.  Its cocycle follows either from
these exact cancellations or by direct multiplication.  Thus it
descends on the original free properly discontinuous toric quotient.
Let $L_0$ be its twelfth power.  On the punctured cusp the exact
ratio $L_0\otimes L_*^{-1}$ has factors $q^{-6n_u}$.

\item \emph{The inverse used on every annulus.}
For a line with normalized lattice factors $1$ on the two real
periods and $e^{2\pi i n_T^tl(t)}$ on the columns of $T$, the form
$A=a(t)^t d\zeta$ is a holomorphic connection when
\begin{equation}\label{inv:eq:atlas-connection-inverse}
 T'^t a=-2\pi i l',\qquad
 (T'^t)^{-1}=
 \frac1{\tau'\beta'-6\mu'^2}
 \begin{pmatrix}\beta'&-\mu'\\-6\mu'&\tau'\end{pmatrix}.
\end{equation}
Indeed translation by $Tn_T$ adds
$a^tT'n_T\,dt=-d\log(e^{2\pi in_T^tl})$.

Every use of this inverse is on one of the three annuli above.
By \cref{inv:A2.1},
$6\mu_z^2-\tau_z\beta_z=\tau_z^2E_4(\tau)^2/E_6(\tau)$,
with the apparent poles removed, has order two at the order-three
point, order zero at the order-four point, and no further zeros on
the regular locus.  A local change of base coordinate multiplies
the determinant by the square of a nonzero derivative.  Thus
$\det T'$ is nonzero at every point of every annular cover used
here.  No inverse is taken at a special fibre.  In the cusp variable
its leading term is $1/(4\pi^2q^2)$, in agreement with the exact
logarithmic period matrix.

\item \emph{The cusp overlap connection.}
For $F_0=L_0\otimes L_*^{-1}$, use
$l=(-6\log q/(2\pi i),0)^t$.  Its connection is
\begin{equation}\label{inv:eq:atlas-cusp-connection}
 A_0=a_0^td\zeta,\qquad
 a_0=\frac6q(T'^t)^{-1}\binom10
 =\frac6{q(\tau'\beta'-6\mu'^2)}\binom{\beta'}{-6\mu'}.
\end{equation}
It is single-valued: a logarithm turn adds
$\operatorname{diag}(1,-1)$ to $T$ and changes neither $T'$ nor
$\zeta$.  In the actual cusp frame $f_{ij}^{12}$ the form is
$A_0+12d\log f_{ij}$.  This transformation retains the full
units \eqref{inv:eq:atlas-cusp-unit}.

\item \emph{The finite overlap connections and all logarithm turns.}
The actual collar map, as in
\cref{wb:thm:retained-period-deformation}, is
\[
 H(s,\zeta)=(s,\zeta+\sigma),\qquad
 \sigma=h\Pi v_j,\qquad h=\frac{\log s}{2\pi i}.
\]
The identities $h(g_js)=h(s)-1/m_j$ and $A_jv_j=v_j$ give the
exact square $Hg_j^{\rm aff}=g_j^{\rm lin}H$.
For $F=H^*L_*\otimes L_j^{-1}$ its lattice factors are
$e^{2\pi i\,12n_T^tD\sigma}$.  Put
$v_q=(v_u,v_\gamma)^t$ and $S=v_q^tD\Pi v_j$.  Then
\begin{equation}\label{inv:eq:atlas-finite-connection}
 A=a^td\zeta-12\pi i\,d(h^2S),\qquad
 a=-24\pi i(T'^t)^{-1}D\sigma'.
\end{equation}
The last differential is a base one-form.
Equation \eqref{inv:eq:atlas-connection-inverse} proves lattice descent.
For an integer logarithm turn $h\mapsto h+k$, the frame changes by
$a_{kv_j}(\zeta+\sigma)^{12}$.  The vertical change in
\eqref{inv:eq:atlas-finite-connection} is
$-24\pi i kDv_q$.  The base change is exactly
\[
 -12\pi i\,d((2kh+k^2)S)
 =-d\log a_{kv_j}(\zeta+\sigma)^{12}
      +24\pi i k v_q^tD\,d\zeta.
\]
Thus the complete connection descends across every logarithm turn,
with no missing base term.

For the remaining finite cyclic descent its actual line multiplier is
\[
 k_F=\frac{e^{2\pi i\,12F_j(\zeta+\sigma)}}{k_j(s,\zeta)}.
\]
Let $k_F^{(r)}$ be the product over the first $r$ actual affine
iterates.  The finite sum
\begin{equation}\label{inv:eq:atlas-averaged-connection}
 A^{\rm inv}=\frac1{m_j}\sum_{r=0}^{m_j-1}
       \bigl((g_j^r)^*A+d\log k_F^{(r)}\bigr)
\end{equation}
is a connection on the same line.  The exact power relation in
\eqref{inv:eq:atlas-affine-lift} makes cyclic permutation of the
summands valid, so it descends to the original annulus.  For
$F_j=L_j\otimes L_*^{-1}$ use the dual connection $-A^{\rm inv}$.
This is the tensor-dual map, including its sign; no line factors are
discarded.  These holomorphic connections need not be flat:
$da$ can supply a $(2,0)$ curvature term.

\item \emph{The holomorphic one-form cocycle.}
Choose the specified line frames in each same-piece chart.  If
the actual line-coordinate map is $v_b=g_{ba}v_a$, put
\begin{equation}\label{inv:eq:atlas-cocycle}
 c_{ab}=\frac{i}{2\pi}d\log g_{ba}\quad\hbox{within a piece},
 \qquad
 c_{ab}=-\frac{i}{2\pi}A_{F_j,(a,b)}\quad
       (a\text{ middle},\ b\text{ in piece }j),
\end{equation}
and $c_{ba}=-c_{ab}$.
On a triple with two middle charts, the ratio frame of $F_j$
changes by the inverse middle transition, and its connection changes
by $+d\log g_*$.  On a triple with two special charts it changes
by the special transition, and its connection changes by
$-d\log g_j$.  These signs give
$c_{ab}+c_{bc}=c_{ac}$ in both cases.  The special pieces do not
intersect one another.  This proves the cocycle identity on every
component of every overlap.

We have glued a cocycle of holomorphic one-forms.  We have not
specified cross-piece functions with those logarithmic derivatives:
the connections are not assumed closed.  In particular no global
line with nonzero fibre Chern class is being glued.

\item \emph{A completely specified smooth partition.}
In every original uniformizing domain enumerate all rational closed
Euclidean balls $B(c_n,R_n)$ compactly contained in that domain.
Use a fixed enumeration of the original charts by piece, cone label
and finite deck word.  Order the balls by integer height bounding
the chart index and all coordinate numerators and denominators,
then lexicographically, removing repeats.  Each height block is
finite.  The half-radius balls cover every chart.

Set
\[
 E(t)=\begin{cases}e^{-1/t}&t>0,\\0&t\le0,\end{cases}
 \quad \chi(t)=\frac{E(t)}{E(t)+E(1-t)},\quad
 b_n(u)=\chi\!\left(
       \frac{R_n^2-\|u-c_n\|^2}{3R_n^2/4}\right).
\]
Extend $b_n$ by zero outside its ball.  It is one on the half-radius
ball and is flat on its support boundary.  For the actual étale map
$\pi_n$, define
\begin{equation}\label{inv:eq:atlas-weights}
 \begin{aligned}
 S_n(x)&=\sum_{u\in\pi_n^{-1}(x)}b_n(u),\qquad
 a_n(x)=\chi(2S_n(x)),\\
 P_1&=1,\qquad P_{n+1}=P_n(1-a_n),\\
 R(s)&=\chi(2s)/s\quad(s>0),\qquad R(0)=0,\\
 w_n(u)&=P_n(\pi_nu)R(S_n(\pi_nu))b_n(u).
 \end{aligned}
\end{equation}
The function $R$ is smooth at zero because its numerator is flat.
Each trace $S_n$ is locally a finite smooth sum: its source support
is compact and the map is étale, hence has discrete fibres.  Every
point has a lift in an open half-radius ball.  For some $N$ one therefore
has $S_N\ge1$ on a neighbourhood, after shrinking that neighbourhood
if necessary.  There $a_N=1$, so all subsequent weights vanish.
Consequently the weights are locally finite.  The finite telescoping
product gives
\[
 \sum_n\sum_{u\in\pi_n^{-1}(x)}w_n(u)=1.
\]
This construction also handles original chart maps that are not
globally injective; it retains all their local inverse branches.

For completeness, every derivative in the final coefficients is
given by the following finite recursion.  In a local inverse branch
write $\mathcal D=\partial/\partial\bar z_l$ and
$v_n=(R_n^2-\|u-c_n\|^2)/(3R_n^2/4)$.  Then
\begin{equation}\label{inv:eq:atlas-weight-derivative}
 \begin{aligned}
 \mathcal D b_n&=-\frac4{3R_n^2}\chi'(v_n)
        \sum_j(u_j-c_{n,j})\mathcal D\bar u_j,\\
 \mathcal D S_n&=\sum_{u\in\pi_n^{-1}(x)}\mathcal D b_n,
 &\mathcal D a_n&=2\chi'(2S_n)\mathcal D S_n,\\
 \mathcal D P_1&=0,
 &\mathcal D P_{n+1}&=(1-a_n)\mathcal D P_n-P_n\mathcal D a_n,\\
 R'(s)&=\frac{2s\chi'(2s)-\chi(2s)}{s^2}\quad(s>0),
 &R'(0)&=0,\\
 \mathcal D w_n&=R(S_n)b_n\mathcal D P_n
       +P_nR'(S_n)b_n\mathcal D S_n+P_nR(S_n)\mathcal D b_n.
 \end{aligned}
\end{equation}
The derivatives $\mathcal D\bar u_j$ are the conjugate Jacobians
of the explicit inverse chart maps in step~1.  Finally
\[
 \chi'(t)=
 \frac{E'(t)E(1-t)+E(t)E'(1-t)}{(E(t)+E(1-t))^2},
 \qquad E'(t)=\begin{cases}e^{-1/t}/t^2&t>0,\\0&t\le0.\end{cases}
\]
Thus no unspecified smooth correction is hidden in the coefficients.

\item \emph{The primitive scalar cocycle.}
For a complex parameter $z$, use lattice factors
$e^{-2\pi iz\gamma(\lambda)}$ and middle linear factors one.
The finite affine factor is
$e^{-2\pi iz\gamma(v_j)/m_j}$; its $m_j$th power is the actual
period factor of $v_j$.  The cusp period factor is
$e^{-2\pi izn_\gamma}$, which extends in every toric chart.
The actual finite collar comparison is
$e^{-2\pi iz\gamma(v_j)h(s)}$.
Its changes under $h\mapsto h-1/m_j$ and $h\mapsto h+k$ are
exactly those required by the affine generator and the period $kv_j$.

Write every resulting transition as $e^{2\pi iz f_{ab}}$.
Its coefficients are respectively
$-\gamma(\lambda)$, $0$, $-\gamma(v_j)/m_j$,
$-n_\gamma$ and $-\gamma(v_j)h(s)$ on those arrows;
pure toric frame changes have coefficient zero.
Composition and inversion use addition and negation.
Since the identities hold for every complex $z$, the $f_{ab}$
satisfy an exact additive cocycle, rather than one only modulo
integers.

\item \emph{The evaluated global forms and their normalization.}
In a chart $a$, with the indicated component of each fibre-product
overlap, put
\begin{equation}\label{inv:eq:atlas-final-forms}
 H_a=\sum_n\sum_u w_n(u)f_{a,n}(x,u),\quad
 B_a=\sum_n\sum_u w_n(u)c_{a,n}(x,u),\quad
 G=\bar\partial H_a,\quad Q=\frac1{12}\bar\partial B_a.
\end{equation}
The cocycles and sum of weights give
$H_a-H_b=f_{ab}$ and $B_a-B_b=c_{ab}$.
Their right sides are holomorphic.  Therefore these formulas agree
on every overlap and define global $\bar\partial$-closed forms.
Writing $c_{a,n}=\sum_kc_{a,n,k}dz_k$, their literal coefficients are
\begin{equation}\label{inv:eq:atlas-final-coefficients}
 G_{\bar l}=\sum_n\sum_u f_{a,n}(x,u)
                       \frac{\partial w_n(u)}{\partial\bar z_l},
 \qquad
 Q_{k\bar l}=-\frac1{12}\sum_n\sum_u c_{a,n,k}(x,u)
                       \frac{\partial w_n(u)}{\partial\bar z_l}.
\end{equation}
Every coefficient is a locally finite sum of the specified period,
Cartier, connection and elementary cutoff expressions.  The minus
sign in the second formula is from moving $d\bar z_l$ past $dz_k$.

On a regular fibre, the one-form cocycle is that of the line with
Chern class $12\eta$, with the Chern sign fixed in step~2.
Consequently $i_b^*[Q]=\eta$.  This is a statement about the exact
Dolbeault quotient; the resulting coefficient form need not be
pointwise translation invariant on the fibre.
For the scalar cocycle, $f_{ab}=-\Delta r_\gamma$, so its fibre
image is $\bar\partial r_\gamma$, explicitly
\[
 \frac{i}{2\Delta}
 \left(d\bar\zeta_2-\frac{\operatorname{Im}\mu}{\operatorname{Im}\tau}
                      d\bar\zeta_1\right).
\]
This is \CompendiumRef{I}{wb:eq:bcpr-gamma-form}; injectivity in
\cref{inv:C3.1} gives $[G]=g_\gamma$.
At $z=1$ the scalar line has the actual finite root characters and
collar frames of $\mathcal O(S_1-S_2)$.  Its twelfth power is
\[
 \mathcal O(12S_1-12S_2)
 =f^*\mathcal O(4p_1-3p_2)\simeq f^*\mathcal O(1),
\]
in agreement with the original sign and factor in \cref{inv:C3.2}.

\item \emph{Why the correction is complex.}
Suppose a real $(1,1)$ form with the stated fibre class were
$\bar\partial$ closed.  Reality would give $\partial Q=0$, hence
$dQ=0$.  Since $H^2(X;\mathbb R)=0$, write $Q=dA$.
On a regular compact fibre, Dolbeault equality with $\eta$ and
Stokes' theorem give $\int_F(i^*Q)^2=\int_F\eta^2=-12$.
But $(i^*Q)^2=d(i^*A\wedge i^*Q)$ has integral zero.
This proves the assertion and shows why the complex connection
correction above cannot be replaced by gluing real curvatures.
\end{enumerate}
\end{proof}
